% !PRINTED 781
\weberpage{781}
%<<CONT
sary. Its value is equal to the square root of the quotient
of the discriminants
\[
\Delta (1, r, r^2, \ldots, r^{\nu-1}) : \Delta (\omega_1, \omega_2, \ldots, \omega_\nu),
\]
in which $\omega_1, \omega_2, \ldots, \omega_\nu$ is a minimal basis (\S.\,144, 145).

From the system of the $p$ we exclude all primes which divide the discriminant of the cyclotomic equation and all which divide $m$ and $a$, whose number is in any case finite, and we denote by $\mathfrak{p}$ a prime ideal dividing $p$. Then the congruence
\[
r^p \equiv r \ (\mathrm{mod}\ \mathfrak{p})
\]
is satisfied always, but also only, when

\weq{16}{p \equiv 1 \ (\mathrm{mod}\ m)}

holds, that is, when $p$ belongs to the exponent 1, and hence when $r^p = r$; and under the same condition we therefore have also, by (15), for every integer $\omega$ in $\Omega_m$:
\[
\omega^p \equiv \omega \ (\mathrm{mod}\ \mathfrak{p}).
\]
But by \S.\,162 this is the necessary and sufficient condition for the prime $p$, which does not divide the field discriminant of $\Omega_m$, to split into nothing but prime ideals of the first degree.

If, then, we denote by
\[
P_f(s) = \prod (1 - p^{-sf})
\]
the product extended over all primes $p$ which belong to the exponent $f$, it follows from \S.\,192 that, if $f>1$, $P_f(1)$ is finite and different from zero, and that

\weq{17}{\dfrac{s-1}{P_1(s)^{\nu}}}

\gesperrt{is, for} $s=1$, \gesperrt{finite and different from zero.}

Now by (13), if $S$ denotes the \gesperrt{finite} product $\Pi(1-p^{-sf})^{-e}$ extended over the \gesperrt{excluded} primes $p$,
\[
Q_1 Q_2 \ldots Q_\mu = S \overset{f}{\Pi} \frac{1}{P_f(s)^{e}},
\]
and since $e = \mu$ for $f = 1$, it follows from (17) that

\weq{18}{(s-1)^{\frac{\mu}{\nu}} Q_1 Q_2 \ldots Q_\mu}

\gesperrt{is, for} $s = 1$, \gesperrt{finite and different from zero.}
